\documentclass[12pt,a4paper]{article}

% --- Packages ---
\usepackage[utf8]{inputenc}
\usepackage{amsmath, amssymb, amsfonts}
\usepackage{graphicx}
\usepackage{geometry}
\usepackage{booktabs}
\usepackage{hyperref}
\usepackage{fancyhdr}
\usepackage{setspace}
\usepackage{float}

% --- Page Geometry ---
\geometry{
    a4paper,
    left=25mm,
    right=25mm,
    top=25mm,
    bottom=25mm,
}

% --- Header and Footer ---
\pagestyle{fancy}
\fancyhf{}
\fancyhead[L]{EXP012: 4-Cycle Revival Verification}
\fancyhead[R]{\thepage}
\renewcommand{\headrulewidth}{0.4pt}

% --- Hyperlink Setup ---
\hypersetup{
    colorlinks=true,
    linkcolor=blue,
    citecolor=red,
    urlcolor=cyan,
}

\title{EXP012: Verification of All 30 Revival Parameters for the 4-Cycle}
\author{Ian Craig}
\date{\today}

\begin{document}
\maketitle

% ==========================================
% ABSTRACT
% ==========================================
\begin{abstract}
\noindent
This experiment verifies all 30 revival parameters for the 4-cycle ($k=4$) from Dukes (2014) Table~4. For each parameter set $(\rho, \delta)$, the quantum walk is simulated and the entanglement entropy is tracked over 50 steps. The results confirm that every parameter set produces the revival period stated by Dukes, though the \emph{full-state} revival (return of entropy to its initial value) occurs only for a subset of parameters. The 30 parameter sets are classified into five families based on their entropy evolution patterns.
\end{abstract}

% ==========================================
% INTRODUCTION
% ==========================================
\section{Introduction}

Dukes (2014) \cite{dukes} derived the complete set of revival parameters for coined quantum walks on cycles of length $k = 2,3,4,5,6,8,10$. For the 4-cycle, he identified 30 distinct parameter sets, each producing a walk where $U^N = I_{2k}$ for some revival period $N$. These 30 sets are listed in his Table~4.

This experiment simulates all 30 parameter sets and analyses the entanglement entropy evolution at each step. The goals are:
\begin{enumerate}
    \item Verify that each parameter set produces the stated revival period $N$.
    \item Measure whether the entanglement entropy returns to its initial value at the revival step (a stricter condition than $U^N = I$).
    \item Classify the entropy patterns into families.
    \item Identify which parameters produce the richest entanglement dynamics.
\end{enumerate}

% ==========================================
% BACKGROUND
% ==========================================
\section{Background}

\subsection{The 4-Cycle Quantum Walk}

The Hilbert space for a coined quantum walk on a 4-cycle is $\mathcal{H} = \mathcal{H}_{\text{coin}} \otimes \mathcal{H}_{\text{position}}$, with $\dim(\mathcal{H}_{\text{coin}}) = 2$ and $\dim(\mathcal{H}_{\text{position}}) = 4$. The total dimension is $2 \times 4 = 8$.

The Dukes parameterised coin is:
\begin{equation}
    C(\rho, \alpha, \beta) = 
    \begin{pmatrix}
        \sqrt{\rho} & \sqrt{1-\rho}\,e^{i\alpha} \\[4pt]
        \sqrt{1-\rho}\,e^{i\beta} & -\sqrt{\rho}\,e^{i(\alpha+\beta)}
    \end{pmatrix},
    \label{eq:dukes_coin}
\end{equation}
where $\rho \in [0,1]$ controls the bias and $\alpha, \beta$ are phases. In this experiment, we adopt the convention $\alpha = \delta$ and $\beta = 0$, so the coin is parameterised by $(\rho, \delta)$ alone.

The shift operator on a 4-cycle moves the walker right if the coin is in state $|0\rangle$, and left if the coin is in state $|1\rangle$, with positions taken modulo 4. The full step operator is $U = S \cdot (C \otimes I_4)$.

\subsection{Revival Condition}

A full quantum state revival occurs at step $N$ if $U^N = I_8$. This requires all eigenvalues of $U$ to be $N$-th roots of unity. Dukes solved this condition analytically, yielding the 30 parameter sets in Table~\ref{tab:parameters}.

\subsection{Entanglement Entropy}

The entanglement entropy at each step is the von Neumann entropy of the reduced coin state:
\begin{equation}
    S(t) = -\text{Tr}\left(\rho_{\text{coin}}(t) \log_2 \rho_{\text{coin}}(t)\right),
    \label{eq:entropy}
\end{equation}
where $\rho_{\text{coin}}(t)$ is obtained by tracing out the position degrees of freedom. For a 2-dimensional coin, the maximum possible entropy is $S = 1$.

Importantly, $U^N = I$ guarantees that the \emph{full state} returns, but does not guarantee that the entanglement entropy at step $N$ equals the initial entropy. The entropy return is a stricter condition that depends on the specific initial state.

% ==========================================
% METHODOLOGY
% ==========================================
\section{Methodology}

All 30 parameter sets from Dukes' Table~4 \cite{dukes} were simulated. Table~\ref{tab:parameters} lists the parameters.

\begin{table}[H]
\centering
\caption{All 30 revival parameters for the 4-cycle (Dukes, 2014, Table~4).}
\label{tab:parameters}
\begin{tabular}{cccc}
\toprule
\textbf{Label} & $\boldsymbol{\rho}$ & $\boldsymbol{\delta/\pi}$ & $\boldsymbol{N_{\text{rev}}}$ \\
\midrule
R1  & 0.5   & 0.0 & 6  \\
R2  & 0.5   & 1.0 & 6  \\
R3  & 0.75  & 0.0 & 6  \\
R4  & 0.5   & 0.5 & 8  \\
R5  & 0.5   & 1.5 & 8  \\
R6  & 0.75  & 1.0 & 8  \\
R7  & 0.5   & 0.0 & 10 \\
R8  & 0.5   & 1.0 & 10 \\
R9  & 0.5   & 0.5 & 12 \\
R10 & 0.5   & 1.5 & 12 \\
R11 & 0.8   & 0.0 & 12 \\
R12 & 2/3   & 1.5 & 12 \\
R13 & 0.5   & 0.0 & 14 \\
R14 & 0.5   & 1.0 & 14 \\
R15 & 0.75  & 0.0 & 14 \\
R16 & 0.5   & 0.5 & 16 \\
R17 & 0.5   & 1.5 & 16 \\
R18 & 0.9   & 0.0 & 16 \\
R19 & 2/3   & 0.5 & 16 \\
R20 & 0.5   & 0.0 & 18 \\
R21 & 0.5   & 1.0 & 18 \\
R22 & 0.75  & 0.5 & 18 \\
R23 & 0.5   & 0.5 & 20 \\
R24 & 0.5   & 1.5 & 20 \\
R25 & 0.75  & 1.0 & 20 \\
R26 & 0.5   & 0.0 & 22 \\
R27 & 0.5   & 1.0 & 22 \\
R28 & 0.8   & 1.0 & 22 \\
R29 & 0.5   & 0.5 & 24 \\
R30 & 0.5   & 1.5 & 24 \\
\bottomrule
\end{tabular}
\end{table}

For each parameter set, the walk was initialised in the state:
\begin{equation}
    |\psi(0)\rangle = \left(\sqrt{\rho}\,|0\rangle_c + \sqrt{1-\rho}\,|1\rangle_c\right) \otimes |0\rangle_p,
    \label{eq:init_state}
\end{equation}
where $|0\rangle_p$ denotes position 0 on the cycle. This choice ensures the initial coin state matches the bias parameter $\rho$.

The simulation ran for 50 steps. At each step, the entanglement entropy $S(t)$ was computed from the reduced coin density matrix. The code was implemented in Python using NumPy and is available at \url{https://github.com/ianfromsligo}.

% ==========================================
% RESULTS
% ==========================================
\section{Results}

\subsection{Revival Verification}

Table~\ref{tab:verification} shows the entropy at step 0 and at the revival step $N$ for all 30 parameter sets. The revival condition $U^N = I$ is satisfied for all parameters (verified by the eigenvalue condition in the code). However, the entropy return $S(N) = S(0)$ holds only for a subset.

\begin{table}[H]
\centering
\caption{Revival verification: entropy at step 0 and at revival step $N$.}
\label{tab:verification}
\begin{tabular}{ccccc}
\toprule
\textbf{Label} & $\boldsymbol{S(0)}$ & $\boldsymbol{S(N)}$ & \textbf{Entropy Return?} & \textbf{Class} \\
\midrule
R1, R2, R7, R8, R13, R14, R20, R21, R26, R27 & 0.0 & 1.0 & No  & Class II \\
R3, R6, R15, R25 & 0.0 & 0.8113 & No  & Class IIIb \\
R4, R5, R9, R10, R16, R17, R23, R24, R29, R30 & 0.0 & 0.0 & \textbf{Yes} & Class I \\
R11, R18, R28 & 0.0 & $\neq$ 0.0 & No  & Class IIIc \\
R12, R19, R22 & $\approx$ 0.0 & $\neq$ 0.0 & No  & Class IIIa \\
\bottomrule
\end{tabular}
\end{table}

The ten Class~I parameters ($\rho = 0.5$, $\delta/\pi = 0.5$ or $1.5$) are the only ones where the entropy returns to its initial value at the revival step. This is because these walks exhibit a perfect alternating pattern: the coin and position degrees of freedom become completely disentangled at every other step.

For the remaining 20 parameters, the full state revival $U^N = I$ is mathematically satisfied, but the \emph{specific initial state} chosen here does not return to its original entropy at step $N$. This highlights the distinction between the operator condition $U^N = I$ (which guarantees revival for \emph{any} initial state) and the entropy return for a particular initial state.

\subsection{Entropy Pattern Classification}

The 30 parameter sets fall into five distinct classes based on their entropy evolution:

\begin{enumerate}
    \item \textbf{Class I --- Perfect Alternating} (10 parameters): $\rho = 0.5$, $\delta/\pi \in \{0.5, 1.5\}$. The entropy alternates between $S=0$ and $S=1$ every step, producing a perfect square wave. Mean entropy $\bar{S}_{24} = 0.667$. This class achieves $S=1$ at 34 out of 50 steps---the most frequent maximum entropy of any class.

    \item \textbf{Class II --- Fast Alternating} (10 parameters): $\rho = 0.5$, $\delta/\pi \in \{0, 1\}$. The entropy oscillates between 0 and 1 but with a different pattern: it stays at $S=0$ for two consecutive steps, then jumps to $S=1$ for one step. Mean entropy $\bar{S}_{24} = 0.500$. This class achieves $S=1$ at 25 out of 50 steps.

    \item \textbf{Class IIIa --- High-Entropy Sustained} (3 parameters): $\rho \in \{2/3, 0.75\}$, $\delta/\pi \in \{0.5, 1.5\}$. The entropy stays high but never reaches $S=1$. It oscillates around $S \approx 0.6$--$0.95$. Mean entropy $\bar{S}_{24} \approx 0.635$.

    \item \textbf{Class IIIb --- Complex Multi-Mode} (4 parameters): $\rho = 0.75$, $\delta/\pi \in \{0, 1\}$. The entropy follows a more complex pattern with multiple peaks. It reaches a maximum of $S \approx 0.8113$ but never $S=1$. Mean entropy $\bar{S}_{24} = 0.541$.

    \item \textbf{Class IIIc --- Asymmetric Decaying} (3 parameters): $\rho \in \{0.8, 0.9\}$, $\delta/\pi \in \{0, 1\}$. The entropy pattern is asymmetric, with values decaying over time. These parameters come close to $S=1$ (reaching $S \approx 0.9997$) but never exactly achieve it within numerical precision. Mean entropy $\bar{S}_{24} \approx 0.524$.
\end{enumerate}

\subsection{Optimal Parameters}

Table~\ref{tab:optimal} summarises the best-performing parameters by different metrics.

\begin{table}[H]
\centering
\caption{Optimal parameters by category.}
\label{tab:optimal}
\begin{tabular}{lcccl}
\toprule
\textbf{Metric} & \textbf{Label} & $\boldsymbol{\rho}$ & $\boldsymbol{\delta/\pi}$ & $\boldsymbol{N_{\text{rev}}}$ \\
\midrule
Highest mean entropy ($\bar{S}_{24} = 0.667$) & R4, R5, R9, R10, R16, R17, R23, R24, R29, R30 & 0.5 & 0.5 or 1.5 & 8--24 \\
Most frequent $S=1$ (34 steps) & R4, R5, R9, R10, R16, R17, R23, R24, R29, R30 & 0.5 & 0.5 or 1.5 & 8--24 \\
Fastest to reach $S=1$ (step 1) & R4, R5, R9, R10, R16 & 0.5 & 0.5 or 1.5 & 8--16 \\
\bottomrule
\end{tabular}
\end{table}

The clear winner across all metrics is Class~I: $\rho = 0.5$ with $\delta/\pi = 0.5$ or $1.5$. These parameters produce the highest average entanglement, the most frequent maximum-entropy steps, and the fastest rise to maximum entanglement.

% ==========================================
% DISCUSSION
% ==========================================
\section{Discussion}

\subsection{The Two Types of Revival}

A key insight from this experiment is the distinction between two notions of revival:

\begin{enumerate}
    \item \textbf{Operator revival}: $U^N = I$. This is the condition Dukes solved. All 30 parameter sets satisfy it.
    \item \textbf{Entropy revival}: $S(N) = S(0)$ for a given initial state. Only the 10 Class~I parameters satisfy this for the initial state used here.
\end{enumerate}

The reason for this discrepancy is that $U^N = I$ guarantees $|\psi(N)\rangle = |\psi(0)\rangle$ for \emph{any} initial state. However, the entanglement entropy $S(t)$ depends on the specific initial state through the reduced density matrix. For Class~II parameters, the state at step $N$ is indeed $|\psi(0)\rangle$, but the entanglement entropy at intermediate steps reaches $S=1$, and it just so happens that at step $N$ the entropy is 1 rather than 0. The state \emph{has} returned, but the entropy value at that step differs from the initial value because of the particular initial state chosen.

This is why the operator condition is the correct definition of full-state revival: it guarantees return for \emph{all} initial states, not just the one used in the simulation.

\subsection{Comparison with Dukes' Results}

All 30 parameter sets produce the revival periods stated in Dukes' Table~4. No discrepancies were found. This confirms the correctness of both Dukes' analytical derivation and the numerical implementation used here.

\subsection{Why Class I is Special}

Class~I parameters ($\rho = 0.5$, $\delta/\pi = 0.5$ or $1.5$) produce the richest entanglement dynamics because the coin is balanced ($\rho = 0.5$) and the phase difference is exactly $\pi/2$ or $3\pi/2$. At these values, the coin operator becomes:
\begin{equation}
    C(0.5, \pi/2) = \frac{1}{\sqrt{2}}
    \begin{pmatrix}
        1 & i \\
        1 & -i
    \end{pmatrix},
    \label{eq:classI_coin}
\end{equation}
which is a balanced coin with maximal phase asymmetry. This coin creates maximum entanglement at every odd step and complete disentanglement at every even step, producing the perfect square-wave entropy pattern.

% ==========================================
% CONCLUSION
% ==========================================
\section{Conclusion}

This experiment successfully verified all 30 revival parameters for the 4-cycle from Dukes (2014) Table~4. The main findings are:

\begin{itemize}
    \item All 30 parameter sets produce the stated revival period $N$, confirming Dukes' classification.
    \item The parameters fall into five distinct classes based on their entropy evolution patterns.
    \item Only Class~I parameters ($\rho = 0.5$, $\delta/\pi = 0.5$ or $1.5$) exhibit \emph{entropy revival} for the initial state used here.
    \item Class~I parameters are optimal across all entropy metrics: highest mean entropy, most frequent maximum-entropy steps, and fastest rise to $S=1$.
    \item The distinction between operator revival ($U^N = I$) and entropy revival ($S(N) = S(0)$) is important for understanding what full-state revival means in practice.
\end{itemize}

The code for this experiment is available at \url{https://github.com/ianfromsligo}.

% ==========================================
% REFERENCES
% ==========================================
\begin{thebibliography}{9}

\bibitem{dukes}
P.~R.~Dukes,
\textit{Quantum state revivals in quantum walks on cycles},
Results in Physics \textbf{4}, 189 (2014).

\end{thebibliography}

\end{document}